% LaTeX companion of 05-CLT.typ. Both formats are editable.
\chapter{Central Limit Theorem}\label{central-limit-theorem}

\section{convergence in distribution}\label{convergence-in-distribution}

\begin{definition}{\kn{convergence in distribution}}

这里的分布采用 \knref{probability distribution} 的定义。我们称一个 seq of random variables \(\left\{ \mathit{X}_{\mathit{n}} \right\}\) converge in distribution to a random variable \(\mathit{X}\), 写作 \(\mathit{X}_{\mathit{n}}\overset{\mathit{d}}{\rightarrow}\mathit{X}\), 如果 \(\mathit{X}\) 的 分布函数 \(\mathit{F}_{\mathit{X}}\) 下所有 右连续的 \(\mathit{x}\) (即 \(\mathit{F}_{\mathit{X}}(\mathit{x}) = \mathit{F}_{\mathit{X}}(\mathit{x} - )\) ), 都有

\[\lim\limits_{\mathit{n}\rightarrow\infty}\mathit{F}_{\mathit{X}_{\mathit{n}}}(\mathit{x}) = \mathit{F}_{\mathit{X}}(\mathit{x})\]

\end{definition}

\begin{qlremark}{Remark}

Convergence in distribution 的意思就是: 随着 \(\mathit{n}\) 越来越大, \(\mathit{X}_{\mathit{n}}\) 的分布是否越来越接近 \(\mathit{X}\) 的分布. 这是一个非常弱的收敛概念, 因为它甚至\textbf{不要求 \(\mathit{X}_{\mathit{n}}\) 和 \(\mathit{X}\) 在同一个概率空间上定义,} 也不要求 \(\mathit{X}_{\mathit{n}}\) 和 \(\mathit{X}\) 之间有任何关系. 只需要 分布函数收敛即可.

并且 notice: 并不要求对所有点 \(\mathit{x}\) 都有 \(\mathit{F}_{\mathit{X}_{\mathit{n}}}(\mathit{x})\rightarrow\mathit{F}_{\mathit{X}}(\mathit{x})\), 而是只要求对 \(\mathit{X}\) 的分布函数 \(\mathit{F}_{\mathit{X}}\) 下所有右连续的点 \(\mathit{x}\). 这是因为如果不加这个限制, 很多直观上应该收敛的情况在数学上就会失败.

比如: 令\(\mathit{X}_{\mathit{n}} = 1/\mathit{n}\), 一个 constant random variable, 那么 \(\mathit{X}_{\mathit{n}}\) 的分布函数 \(\mathit{F}_{\mathit{X}_{\mathit{n}}}\) 是一个 step function. 但是在 \(\mathit{x} = 0\) 处, \(\mathit{F}_{\mathit{X}_{\mathit{n}}}(0) = 0\) for all \(\mathit{n}\), 因而 \(\lim_{\mathit{n}\rightarrow\infty}\mathit{F}_{\mathit{X}_{\mathit{n}}}(0) = 0\).

而 \(\mathit{F}_{\mathit{X}}(0) = 1\).

所以对于这种跳跃不连续点, 如果不跳过, 那么就会得出 \(\mathit{X}_{\mathit{n}}\) 不收敛于 \(\mathit{X}\) 的错误结论. 但是我们知道, 分布函数作为一个单调递增的函数, 是 a.e. continuous 的, 因而这些点是可以忽略掉的.

\end{qlremark}

\section{Characterization of a distribution}\label{characterization-of-a-distribution}

\subsection{moment generating function}\label{moment-generating-function}

\begin{definition}{\kn{moment generating function}}

对于一个随机变量 \(\mathit{X}\), 其 moment generating function (MGF) 定义为

\[\mathit{M}_{\mathit{X}}(\mathit{t}) = \mathbb{E}\lbrack\mathit{e}^{\mathit{t}\mathit{X}}\rbrack,\quad\mathit{t} \in \mathbb{R}\]

\end{definition}

\begin{qlremark}{Remark}

为什么这个东西叫做 moment generating function 呢? 因为如果 \(\mathit{M}_{\mathit{X}}(\mathit{t})\) 在 \(\mathit{t} = 0\) 的某个 neighborhood 内存在, 那么我们可以通过对 \(\mathit{M}_{\mathit{X}}(\mathit{t})\) 求导来得到 \(\mathit{X}\) 的各阶矩:

\end{qlremark}

\begin{proposition}

如果 \(\mathit{M}_{\mathit{X}}(\mathit{t})\) 在 \(\mathit{t} = 0\) 的某个 neighborhood 内存在, 则 \(\mathit{X}\) 的 \(\mathit{n}\) 阶矩可以表示为

\[\mathbb{E}\lbrack\mathit{X}^{\mathit{n}}\rbrack = \mathit{M}_{\mathit{X}}^{(\mathit{n})}(0)\]

其中 \(\mathit{M}_{\mathit{X}}^{(\mathit{n})}(0)\) 表示 \(\mathit{M}_{\mathit{X}}(\mathit{t})\) 在 \(\mathit{t} = 0\) 处的 \(\mathit{n}\) 阶导数.

\end{proposition}

\subsection{characteristic function}\label{characteristic-function}

\section{Central Limit Theorem}\label{central-limit-theorem-1}

\begin{qlremark}{Remark}

\knref{De Moivre--Laplace theorem} 给出了二项分布到正态分布的先导特例。

\end{qlremark}

\begin{theorem}{\kn{Lindeberg-Levy Central Limit Theorem}}

对于任意一个 seq of i.i.d. random variables \(\left\{ \mathit{X}_{\mathit{i}} \right\}\) with mean \(\mathit{\mu}\) and variance \(\mathit{\sigma}^{2} < \infty\), set \(\mathit{S}_{\mathit{n}} = \mathit{X}_{1} + \mathit{X}_{2} + \cdots + \mathit{X}_{\mathit{n}}\) for each \(\mathit{n}\).

我们有:

\[\frac{\mathit{S}_{\mathit{n}} - \mathit{n}\mathit{\mu}}{\sqrt{\mathit{n}\mathit{\sigma}^{2}}}\overset{\mathit{d}}{\rightarrow}\mathit{N}(0,1)\]

\end{theorem}

\begin{qlremark}{Remark}

LLN 告诉我们, 当 \(\mathit{n}\) 越来越大时, \(\mathit{S}_{\mathit{n}}/\mathit{n}\) 越来越接近 \(\mathit{\mu}\). 每个 \(\mathit{X}_{\mathit{i}}\) 随机取样, 不论结果如何, 其平均值大概率都是 \(\mathit{\mu}\).\\
CLT 则是进一步告诉我们这个接近的过程是如何发生的: 随着 \(\mathit{n}\) 越来越大, \(\mathit{S}_{\mathit{n}}\) 的分布会越来越接近一个正态分布, 且波动的幅度(标准差)数量级为 \(\sqrt{\mathit{n}}\). \textbf{因而平均值 \(\mathit{S}_{\mathit{n}}/\mathit{n}\) 的分布会越来越接近一个均值为 \(\mathit{\mu}\), 标准差为 \(\mathit{\sigma}/\sqrt{\mathit{n}}\) 的正态分布}

\[\frac{\mathit{S}_{\mathit{n}}}{\mathit{n}} = \mathit{\mu} + \frac{\mathit{\sigma}}{\sqrt{\mathit{n}}}\mathit{Z},\quad\mathit{Z} \sim \mathit{N}(0,1)\]

\textbf{随着\(\mathit{n}\) 越来越大, 这个标准差会越来越小, 因而 其极限坍缩为一个 constant \(\mathit{\mu}\).}\\

\end{qlremark}

\begin{qlremark}{Remark}

CLT 的一个主要的应用价值, 就是可以用来估计 由多个 i.i.d. 随机因素共同作用的现象的分布.

统计学中常用于对 confidence interval 的估计, 以及 hypothesis testing, 通过计算如 \(\mathbb{P}\left( {\frac{\mathit{S}_{\mathit{n}} - \mathit{n}\mathit{\mu}}{\sqrt{\mathit{n}\mathit{\sigma}^{2}}} \in \mathit{A}} \right)\) 这种形式的概率来进行推断.\\

\end{qlremark}

在进行证明前, 我们先看一些 applications of CLT. 可能会帮助我们更好地理解 CLT 的意义.

\subsection{applications of CLT}\label{applications-of-clt}

\begin{example}[(判断 coin 是否 fair)]

我们有两个 coins, 想要判断它们是否是 fair coin. 我们可以 toss 这个 coin \(\mathit{n}\) 次, 记录下每次 toss 的结果, 记为 \(\mathit{X}_{1},\mathit{X}_{2},\cdots,\mathit{X}_{\mathit{n}}\), 其中 \(\mathit{X}_{\mathit{i}} = 1\) if the \(\mathit{i}\)-th toss is heads.

现在: 观测到第一个 coin 100 次 toss 中有 38 次是 heads, 第二个 coin 100 次 toss 中有 43 次是 heads.

\begin{qlsolution}{Solution}

假设这两个 coins 是 fair coin. 那么 \(\mathit{X}_{\mathit{i}}\) 是 i.i.d. Bernoulli random variables with parameter \(\mathit{p} = 0.5\), 因而 \(\mathit{\mu} = \mathit{p} = \frac{1}{2}\), \(\mathit{\sigma}^{2} = \mathit{p}(1 - \mathit{p}) = \frac{1}{4}\). 从而根据 CLT,

\[\begin{aligned}
\mathbb{P}(\mathit{S}_{100} < 38) &= \mathbb{P}\left( {\frac{\mathit{S}_{100} - 100 \cdot \frac{1}{2}}{\sqrt{100 \cdot \frac{1}{4}}} < \frac{38 - 100 \cdot \frac{1}{2}}{\sqrt{100 \cdot \frac{1}{4}}}} \right) \\
 &\approx \mathbb{P}\left( {\mathit{Z} < \frac{38 - 50}{5}} \right) \\
 &\approx \mathbb{P}(\mathit{Z} < - 2.4) \approx 0.0082 < 0.01
\end{aligned}\]

因而这个第一个 coin 很可能不是 fair coin. 同样的方法计算出 \(\mathbb{P}\left( {\mathit{S}_{100} \leq 43} \right) \approx 0.0887\), 因而第二个 coin 虽然也可疑不是 fair coin, 但是不如第一个 coin 可疑. 如果以 0.05 作为显著性水平, 那么我们可以拒绝第一个 coin 是 fair coin 的假设.\\

\end{qlsolution}

\end{example}

\begin{example}[(样本量需求的计算)]

工厂生产了一批电线, 我们想知道它们的平均断裂强度 \(\mathit{\mu}\) 是多少. 遂抽取 \(\mathit{n}\) 根电线进行测量, 得到 \(\mathit{X}_{1},...,\mathit{X}_{\mathit{n}}\), 然后计算它们的样本平均值 \({\bar{\mathit{X}}}_{\mathit{n}}\) 来估计 \(\mathit{\mu}\).

已知量: 强度的方差 \(\mathit{\sigma}^{2} = 1/10\).; 我们想估计的是 \(\mathit{\mu}\) 的值. 并且我们希望我们的估计是准的, in the sense that: 误差 \(|{\bar{\mathit{X}}}_{\mathit{n}} - \mathit{\mu}|\) 不超过 0.01 的概率至少为 0.95.

\begin{qlsolution}{Solution}

我们要达到: \(\left. \mathbb{P}( \middle| {\bar{\mathit{X}}}_{\mathit{n}} - \mathit{\mu} \middle| \leq 1/100) \right.\).

我们要把它转成标准的 normal distribution 的形式, 即 \(\mathit{Z} = \frac{\mathit{S}_{\mathit{n}} - \mathit{n}\mathit{\mu}}{\sqrt{\mathit{n}\mathit{\sigma}^{2}}}\) 的形式.

于是我们把 \({\bar{\mathit{X}}}_{\mathit{n}} - \mathit{\mu}\) 写为 \(\frac{\mathit{S}_{\mathit{n}}}{\mathit{n}} - \mathit{\mu}\), 然后两边同时乘 \(\frac{\sqrt{\mathit{n}}}{\mathit{\sigma}}\). 右边的常数项也做相同的变换: \(\frac{1/100}{\mathit{\sigma}/\sqrt{\mathit{n}}} = \frac{\sqrt{\mathit{n}}}{100\mathit{\sigma}}\).

于是原式变为:

\[\left. \mathbb{P}( \middle| \mathit{Z} \middle| \leq \frac{\sqrt{\mathit{n}}}{100\mathit{\sigma}}) \approx 0.95 \right.\]

对于正态分布我们知道

\[\left. \mathbb{P}( \middle| \mathit{Z} \middle| \leq \mathit{x}) = 2\Phi(\mathit{x}) - 1 \right.\]

于是想要:

\[2\Phi\left( \frac{\sqrt{\mathit{n}}}{100\mathit{\sigma}} \right) - 1 \geq 0.95\Longrightarrow\Phi\left( \frac{\sqrt{\mathit{n}}}{100\mathit{\sigma}} \right) \geq 0.975\]

通过查表我们知道当 \(\Phi(\mathit{z}) = 0.975\) 时, \(\mathit{z} \approx 1.96\). 因此要求 \(\frac{\sqrt{\mathit{n}}}{100\mathit{\sigma}} \geq 1.96\), 解得至少需要 \(\mathit{n} \geq (61.98)^{2} \approx 384.16\), floor 一下得到 \(\mathit{n} \geq 385\).

\end{qlsolution}

\begin{qlremark}{Remark}

这里的应用正是我们开头说的: LLN 告诉我们, 当 \(\mathit{n}\) 越来越大时, \(\mathit{S}_{\mathit{n}}/\mathit{n}\) 越来越接近 \(\mathit{\mu}\). 每个 \(\mathit{X}_{\mathit{i}}\) 随机取样, 不论结果如何, 其平均值大概率都是 \(\mathit{\mu}\); 而 CLT 则是进一步告诉我们这个接近的过程是如何发生的:

\[\frac{\mathit{S}_{\mathit{n}}}{\mathit{n}} = \mathit{\mu} + \frac{\mathit{\sigma}}{\sqrt{\mathit{n}}}\mathit{Z},\quad\mathit{Z} \sim \mathit{N}(0,1)\]

随着\(\mathit{n}\) 越来越大, 这个标准差会越来越小, 因而 其极限坍缩为一个 constant \(\mathit{\mu}\). \textbf{因而, 我们可以通过 CLT 来得到: 我们 要用样本均值来估计分布的真实均值的话, 在一定的置信水平下, 需要多少样本量 \(\mathit{n}\) 才能保证我们的估计是准的.}

\end{qlremark}

\end{example}

\subsection{Berry-Esseen Theorem: CLT 的收敛速度}\label{berry-esseen-theorem-clt-ux7684ux6536ux655bux901fux5ea6}

之前的例子中, 我们都 使用了 \(\approx\) 表示: 我们直接把此时的分布近似当作了一个 normal distribution, 来计算一些概率.

但是: 这两个分布的近似行为的本身有多么精准?

下面有一个 theorem 刻画了这件事.

\begin{theorem}{\kn{Berry-Esseen Theorem}}

给定一个 seq of i.i.d. random variables \(\left\{ \mathit{X}_{\mathit{i}} \right\}\) with mean \(\mathit{\mu}\) and variance \(\mathit{\sigma}^{2} < \infty\), 以及 \(\left. \mathbb{E}\lbrack \middle| \mathit{X}_{\mathit{i}} - \mathit{\mu} \middle| {}_{3}\rbrack = \mathit{\rho} < \infty \right.\), set \(\mathit{S}_{\mathit{n}} = \mathit{X}_{1} + \mathit{X}_{2} + \cdots + \mathit{X}_{\mathit{n}}\), 我们有: 对于任意 \(\mathit{x} \in \mathbb{R}\),

\[\left| {\mathbb{P}\left( {\frac{\mathit{S}_{\mathit{n}} - \mathit{n}\mathit{\mu}}{\sqrt{\mathit{n}\mathit{\sigma}^{2}}} \leq \mathit{x}} \right) - \Phi(\mathit{x})} \right| \leq \frac{3\mathit{\rho}}{\mathit{\sigma}^{2}\sqrt{\mathit{n}}}\]

\end{theorem}

\begin{qlremark}{Remark}

给定一个 number of samples \(\mathit{n}\), 我们可以给出此时的 真实概率与正态近似值之间的差值的 一个 upper bound.\\
并且, 这个 bound 不是概率意义上的 bound, 而\textbf{是一个 deterministic bound, 没有任何的随机性}; 并且是 \textbf{uniform across all \(\mathit{x}\) 的 bound.}\\
这是因为这里本身就是两个概率的值之间的差, 而不是样本之间的差. 因而其实这是一个非常强大的 结论, 表面\textbf{我们用正态分布近似估计出的概率 和真实的概率之间的差值是 determinstically controllable 意义上非常小的.}\\
我们马上可以想到: 对于 控制样本数量来达到某个置信水平的问题, 这是一个非常有用的工具, 因为它告诉我们, 当 \(\mathit{n}\) 足够大时, 我们用正态分布近似估计出的概率和真实的概率之间的 差值是非常小并且可以控制的.\\
并且, 这个控制应该是比较精准的. 因为相比其他的控制方法: \textbf{Chebyshev 只需要 first, second order moments 的信息; Hoeffding 只需要有界性的假设, 甚至不需要知道 分布的任何 moment 信息; 而 Berry-Esseen 需要 third order moment 的信息, 因此它应该能够给一个更精准的控制.}

\end{qlremark}

\begin{proof}

//TODO:

\end{proof}

\begin{example}

一个工厂生产电子元件, 每个原件有概率是有缺陷的.

令

\[\mathit{X}_{\mathit{i}} = \left\{ \begin{matrix}
{1,} & {\text{if the}\ \mathit{i}\ \text{th tested component is defective},} \\
{0,} & {\text{otherwise}\ .}
\end{matrix} \right.\]

并假设 i.i.d. with parameter

\[\mathbb{P}(\mathit{X}_{\mathit{i}} = 1) = \mathit{p}\]

其中 \(\mathit{p}\) 是一个未知的参数.

工厂方面表示这个 process 是 under control 的, 并给出了假设: \(\mathit{H}_{0}:\mathit{p} \leq 0.02\).

为了验证这个假设, 一个 quality control manager 需要从生产线上需要随机 \(\mathit{n}\) 个元件进行测试, 并估计出 \(\mathit{p}\) 的值 by 样本均值:

\[{\widehat{\mathit{p}}}_{\mathit{n}} := \frac{1}{\mathit{n}}\sum\limits_{\mathit{i} = 1}^{\mathit{n}}\mathit{X}_{\mathit{i}}\]

我们希望这个估计的误差最多为 0.005, 并且这个误差的置信度至少为 0.99, 即我们希望:

\[\mathbb{P}\left( {\left| {{\widehat{\mathit{p}}}_{\mathit{n}} - \mathit{p}} \right| > 0.005} \right) \leq 0.01\]

那么我们至少需要多少样本量 \(\mathit{n}\) 来达到这个要求呢?

\end{example}

\begin{qlsolution}{Solution}

首先计算样本均值 \({\widehat{\mathit{p}}}_{\mathit{n}}\) 的 mean 和 variance:

\[\mathbb{E}\left\lbrack {\widehat{\mathit{p}}}_{\mathit{n}} \right\rbrack = \mathit{p},\quad\text{Var}\ \left( {\widehat{\mathit{p}}}_{\mathit{n}} \right) = \frac{1}{\mathit{n}^{2}}\sum\limits_{\mathit{i} = 1}^{\mathit{n}}\text{Var}\ \left( \mathit{X}_{\mathit{i}} \right) = \frac{\mathit{p}(1 - \mathit{p})}{\mathit{n}}\]

此时我们有三个办法:

\begin{itemize}
\item
  办法1: \textbf{Chebyshev's inequality}.

  \[\mathbb{P}\left( {\left| {{\widehat{\mathit{p}}}_{\mathit{n}} - \mathit{p}} \right| \geq \mathit{\varepsilon}} \right) \leq \frac{\text{Var}\ \left( {\widehat{\mathit{p}}}_{\mathit{n}} \right)}{\mathit{\varepsilon}^{2}} = \frac{\mathit{p}(1 - \mathit{p})}{\mathit{n}\mathit{\varepsilon}^{2}}\]

  由于 \(\mathit{p}(1 - \mathit{p}) \leq 1/4\) for any \(\mathit{p} \in \lbrack 0,1\rbrack\), 式子可以进一步控制为

  \[\mathbb{P}\left( {\left| {{\widehat{\mathit{p}}}_{\mathit{n}} - \mathit{p}} \right| \geq \mathit{\varepsilon}} \right) \leq \frac{1}{4\mathit{n}\mathit{\varepsilon}^{2}}\]

  我们希望此概率不超过 0.01, 并且~\(\mathit{\varepsilon} = 0.005\), 那么进一步得到需要

  \[\frac{1}{4\mathit{n}(0.005)^{2}} \leq 0.01\]

  解得至少需要 \(\mathit{n} \geq 1000000\).
\item
  办法2: \textbf{Hoeffding's inequality}.\\
  由于 \(0 \leq \mathit{X}_{\mathit{i}} \leq 1\), 我们可以直接使用 Hoeffding's inequality 来控制:

  \[\mathbb{P}\left( {\left| {{\widehat{\mathit{p}}}_{\mathit{n}} - \mathit{p}} \right| \geq \mathit{\varepsilon}} \right) \leq 2\mathit{e}^{- 2\mathit{n}\mathit{\varepsilon}^{2}}\]

  我们 require 了 \(2\mathit{e}^{- 2\mathit{n}(0.005)^{2}} \leq 0.01\), 因而可以解得

  \[\mathit{n} \geq \frac{|\ln(0.005)|}{2(0.005)^{2}} \approx 106,000\]
\item
  办法3: \textbf{Berry-Esseen Theorem}.\\
  首先我们假设报告的 \(\mathit{p}\) 的值 0.02 是正确的, 从而可以计算出:

  \[\mathbb{E}\left\lbrack \mathit{X}_{1} \right\rbrack = \mathit{p},\quad\mathit{\sigma}^{2} := \text{Var}\ \left( \mathit{X}_{1} \right) = \mathit{p}(1 - \mathit{p}) = 0.0196,\quad\mathit{\sigma} = 0.14\]

  并且, 由于 \(\left| {\mathit{X}_{1} - \mathit{p}} \right| \leq 1\), 我们知道偏度 \(\mathit{\rho} := \mathbb{E}\left\lbrack \left| {\mathit{X}_{1} - \mathit{p}} \right|^{3} \right\rbrack < \infty\).

  因而可以应用 Berry-Esseen Theorem. 我们希望:

  \[\mathbb{P}\left( {\left| {{\widehat{\mathit{p}}}_{\mathit{n}} - \mathit{p}} \right| \leq \mathit{\varepsilon}} \right) = \mathbb{P}\left( {\left| {\mathit{S}_{\mathit{n}} - \mathit{n}\mathit{p}} \right| \leq \mathit{n}\mathit{\varepsilon}} \right) \geq 0.99,\quad\mathit{\varepsilon} = 0.005\]

  等价于

  \[\mathbb{P}\left( {\left| \frac{\mathit{S}_{\mathit{n}} - \mathit{n}\mathit{p}}{\mathit{\sigma}\sqrt{\mathit{n}}} \right| \leq \frac{\mathit{\varepsilon}\sqrt{\mathit{n}}}{\mathit{\sigma}}} \right) \geq 0.99\]

  Berry-Esseen Theorem 可以得到: 对于任意的 \(\mathit{x} > 0\), 有

  \[\mathbb{P}\left( {\left| \frac{\mathit{S}_{\mathit{n}} - \mathit{n}\mathit{p}}{\mathit{\sigma}\sqrt{\mathit{n}}} \right| \leq \mathit{x}} \right) \geq 2\Phi(\mathit{x}) - 1 - \frac{6\mathit{\rho}}{\mathit{\sigma}^{2}\sqrt{\mathit{n}}}\]

  因而我们需要选择 \(\mathit{n}\) s.t.

  \[2\Phi\left( \frac{\mathit{\varepsilon}\sqrt{\mathit{n}}}{\mathit{\sigma}} \right) - 1 - \frac{6\mathit{\rho}}{\mathit{\sigma}^{2}\sqrt{\mathit{n}}} \geq 0.99\]

  忽略掉 \(\frac{6\mathit{\rho}}{\mathit{\sigma}^{2}\sqrt{\mathit{n}}}\) 这个很小的项 (计算可以再精细地选择 \(\mathit{n}\)), 我们也可以得到:

  \[\frac{\mathit{\varepsilon}\sqrt{\mathit{n}}}{\mathit{\sigma}} \geq \mathit{z}_{0.995} \approx 2.576\]

  因而至少需要 \(\mathit{n} \geq 5200\).
\end{itemize}

我们可以看到

\begin{itemize}
\item
  通过 Chebyshev's inequality 来控制, 我们需要的样本量是 100 万级别的;
\item
  通过 Hoeffding's inequality 来控制, 我们需要的样本量是 10 万级别的;
\item
  通过 Berry-Esseen Theorem 来控制, 我们需要的样本量是 5000 级别的, 显然更为高效.
\end{itemize}

不过, 这个结果是基于我们假设 \(\mathit{p} = 0.02\) 的前提下得到的. Chebyshev 的逻辑 (保守估计): 我不相信厂家的任何话, 我要建立一个无论真实 \(\mathit{p}\) 是多少都绝对成立的置信区间. 所以我必须用最差的情况 \(\mathit{p} = 0.5\) 来算 \(\mathit{n}\). CLT/Berry-Esseen 的逻辑 (假设检验): 厂家宣称 \(\mathit{p} \leq 0.02\). 在统计学中, 我们通常先假设这个宣称是正确的 (即 \(\mathit{H}_{0}\) 假设). 如果我们用在这个假设下得出的需要的样本量 \(\mathit{n} = 5200\) 去测, 发现结果远超 \(0.02\), 那我们就直接推翻厂家的说法.

\end{qlsolution}

\subsection{proof of CLT}\label{proof-of-clt}
